\documentclass[sigconf,screen]{acmart}

\begin{document}

\title{rocq-piler: A Content-Addressed Proof Oracle for Discharging Separation-Logic Obligations}

\author{Gavin Mendel-Gleason}
\affiliation{%
  \institution{Scidonia}
  \city{Dublin}
  \country{Ireland}
}
\email{gavin@scidonia.com}

\maketitle

\section*{Demo Abstract}

We demonstrate \textbf{rocq-piler}, a Model Context Protocol (MCP) server that
connects Large Language Models (LLMs) to the Rocq proof assistant (formerly Coq)
through its Language Server Protocol and Pétanque execution backends. \textbf{rocq-piler}
is built around \emph{content-addressed proof obligations}: every open goal is
identified by a stable hash of its statement and hypotheses, so an autonomous
agent targets, batches, and revisits obligations by hash rather than by fragile
line numbers or bullet positions. Our motivating use is as a \emph{proof oracle}
for \textbf{axiomander}, a separate open-source deductive verifier that emits
Iris separation-logic obligations; automation discharges most obligations, but a
residual fraction---roughly one in five---requires an expert, and \textbf{rocq-piler}
is the tunable oracle we use to play that role.

\paragraph{What We Will Demonstrate}

\begin{enumerate}
\item \textbf{The hash-driven loop}: \texttt{focus\_proof} enumerates open
obligations with their hashes, hypotheses, and goals; \texttt{insert\_tactics}
closes them by hash; the cycle repeats to \texttt{Qed}, robustly under the
agent's own edits.

\item \textbf{Batch moves}: \texttt{stratify} case-splits a proof and
auto-closes the routine majority with a tactic portfolio, leaving only genuine
survivors as hash-addressable admits; \texttt{close\_admits} sweeps up survivors
speculatively, never regressing the script.

\item \textbf{Speculative execution}: testing tactics with \texttt{dry\_run} and
the Pétanque state API before committing any edit.

\item \textbf{Incremental caching}: re-verifying only changed proofs and their
affected callers, keeping repeated oracle calls cheap.

\item \textbf{Reproducible case study}: an autonomous proof of type preservation
for PCF with mutable references (21 inductive cases, 9 helper lemmas), closed by
two different LLMs through the same MCP interface for a few cents per run---a
faithful stand-in for the hard tail of axiomander's obligations.
\end{enumerate}

\paragraph{Technical Innovation}

Unlike LLM+prover integrations that focus on model training or bespoke
environments (e.g.\ LeanDojo, Baldur), rocq-piler standardises the
\emph{interaction protocol} and contributes a content-addressed obligation model
plus batch, speculative, non-regressing closing moves. Any MCP-capable model can
serve as the oracle, and the tool surface, prompting, and closing portfolios are
all tunable parameters.

\paragraph{Availability}

\begin{itemize}
\item GitHub: \url{https://github.com/scidonia/rocq-piler}
\item License: MIT; published on npm as \texttt{rocq-piler}
\item Requirements: Node.js 20+, Rocq/Coq 8.19+, rocq-lsp
\end{itemize}

\paragraph{Demo Format}

\begin{itemize}
\item Live session: an LLM closing a multi-case induction via
\texttt{stratify}/\texttt{close\_admits} in real time.
\item Audience-suggested goals dispatched through the hash-driven loop.
\item A deep-dive on the content-addressed model and the axiomander oracle
integration.
\end{itemize}

\end{document}
